% Laget automatisk av elevark/lag_elevark.py fra sim/ladd-kule.html. Endringer her blir overskrevet.
\documentclass[11pt]{article}
\usepackage{elevark}
\begin{document}
\elevarkhode{Fysikk 2 · Elektriske felt og krefter}{Ladd kule mellom plater}{https://lektorbo.no/sim/ladd-kule.html}{Simuleringen}

En lett, ladd kule henger i en snor midt mellom to parallelle metallplater. Når det er spenning mellom platene, blir det et homogent elektrisk felt mellom dem, og kula blir trukket til siden. Hvor stor blir vinkelen?

\section*{Oppgaver}

Start med $U = 5{,}0\,\mathrm{kV}$, $d = 10\,\mathrm{cm}$, $q = +10\,\mathrm{nC}$ og $m = 0{,}20\,\mathrm{g}$.

\begin{oppgaveliste}
\oppgave \fase{gjett}{Gjett}Hva skjer med vinkelen hvis du snur spenningen, slik at høyre plate blir positiv? Hva skjer hvis du i stedet gjør kula negativ? Test begge.
\svar
\oppgave \fase{test}{Regn}Regn ut feltstyrken mellom platene, den elektriske kraften på kula og tyngden. Bruk kraftdiagrammet til å vise at $\tan  θ = F_{\mathrm{e}} / G$, og regn ut vinkelen. Sjekk med simuleringen.
\svar
\oppgave \fase{gjett}{Gjett}Hva skjer med vinkelen hvis du dobler spenningen? Blir vinkelen dobbelt så stor? Test med 10~kV.
\svar
\oppgave \fase{forklar}{Forklar}Sett spenningen tilbake til 5,0~kV, og øk avstanden mellom platene fra 10~cm til 20~cm. Hvorfor blir vinkelen mindre? Hva må du gjøre med spenningen for at vinkelen skal bli som før?
\svar
\oppgave \fase{test}{Regn}Kula har massen 0,30~g og ladningen 12~nC, og avstanden er 10~cm. Hvor stor må spenningen være for at snora skal danne 20° med loddlinjen? Test svaret.
\svar
\oppgave \fase{test}{Mål}Velg «Ukjent (regn den ut)». Les av vinkelen, og regn ut ladningen til kula. Er den positiv eller negativ? Sjekk med «Vis ladningen». Prøv flere ukjente ladninger.
\svar
\oppgave \fase{utfordring}{Utfordring}Feltet mellom platene er homogent. Hva betyr det for kraften på kula om den henger nær den ene plata i stedet for midt mellom? Hvorfor er dette ikke tilfellet for kula i simuleringen Elektrisk felt\footnote{\url{lektorbo.no/sim/elektrisk-felt.html}}?
\svar
\end{oppgaveliste}

\end{document}
